% ------------------------------------------------------------------- %
%  Daftar Singkatan, Istilah, dan Lambang
%  AUTHOR: @NoHaitch IF'22
%  version: FINAL
% ------------------------------------------------------------------- %

\chapter*{Daftar Singkatan, Istilah, dan Lambang}
\addcontentsline{toc}{chapter}{\uppercase{Daftar Singkatan, Istilah, dan Lambang}}
\label{daftar-singkatan}

\todo{

\begin{xltabular}{\textwidth}{|l|X|c|}
    \caption*{Daftar singkatan yang digunakan dalam naskah.} \\
    \hline
    \rowcolor[HTML]{EFEFEF}
    \multicolumn{1}{|c|}{\textbf{Singkatan}} & \multicolumn{1}{c|}{\textbf{Penjelasan}} & \makecell{\textbf{Halaman} \\ \textbf{Pertama}} \\
    \hline
    \endfirsthead
    \multicolumn{3}{c}{{Daftar singkatan (lanjutan).}} \\
    \hline
    \rowcolor[HTML]{EFEFEF}
    \multicolumn{1}{|c|}{\textbf{Singkatan}} & \multicolumn{1}{c|}{\textbf{Penjelasan}} & \makecell{\textbf{Halaman} \\ \textbf{Pertama}} \\
    \hline
    \endhead
    \hline
    \multicolumn{3}{r}{\textit{Bersambung ke halaman berikutnya}} \\
    \endfoot
    \hline
    \endlastfoot
    \textbf{FXX} &
    Kode \textit{Fungsional} untuk kebutuhan fungsional sistem, dengan \texttt{XX} menyatakan nomor urut (01--20) &
    \pageref{first-FXX} \\ \hline

    \textbf{NFXX} &
    Kode kebutuhan \textit{Non-Fungsional}, khususnya properti anonimitas, dengan \texttt{XX} menyatakan nomor urut (01--04) &
    \pageref{first-NFXX} \\ \hline
\end{xltabular}

\begin{xltabular}{\textwidth}{|>{\hsize=0.45\hsize}X|>{\hsize=1.55\hsize}X|c|}
    \caption*{Daftar istilah yang digunakan dalam naskah.} \\
    \hline
    \rowcolor[HTML]{EFEFEF}
    \multicolumn{1}{|c|}{\textbf{Istilah}} & \multicolumn{1}{c|}{\textbf{Penjelasan}} & \makecell{\textbf{Halaman} \\ \textbf{Pertama}} \\
    \hline
    \endfirsthead
    \multicolumn{3}{c}{{Daftar istilah (lanjutan).}} \\
    \hline
    \rowcolor[HTML]{EFEFEF}
    \multicolumn{1}{|c|}{\textbf{Istilah}} & \multicolumn{1}{c|}{\textbf{Penjelasan}} & \makecell{\textbf{Halaman} \\ \textbf{Pertama}} \\
    \hline
    \endhead
    \hline
    \multicolumn{3}{r}{\textit{Bersambung ke halaman berikutnya}} \\
    \endfoot
    \hline
    \endlastfoot

    % ISI DI SINI
    \textbf{\textit{Smart contract}} &
    Program yang berjalan pada platform DLT dan mengeksekusi logika bisnis secara otomatis sesuai kondisi yang diprogram &
    \pageref{first-smart-contract} \\ \hline
\end{xltabular}

\begin{xltabular}{\textwidth}{|l|X|c|}
    \caption*{Daftar lambang dan notasi yang digunakan dalam naskah.} \\
    \hline
    \rowcolor[HTML]{EFEFEF}
    \multicolumn{1}{|c|}{\textbf{Lambang}} &
    \multicolumn{1}{c|}{\textbf{Penjelasan}} &
    \makecell{\textbf{Halaman} \\ \textbf{Pertama}} \\
    \hline
    \endfirsthead
    \multicolumn{3}{c}{{Daftar lambang dan notasi (lanjutan).}} \\
    \hline
    \rowcolor[HTML]{EFEFEF}
    \multicolumn{1}{|c|}{\textbf{Lambang}} &
    \multicolumn{1}{c|}{\textbf{Penjelasan}} &
    \makecell{\textbf{Halaman} \\ \textbf{Pertama}} \\
    \hline
    \endhead
    \hline
    \multicolumn{3}{r}{\textit{Bersambung ke halaman berikutnya}} \\
    \endfoot
    \hline
    \endlastfoot

    % ISI DI SINI
    \textbf{$d$} &
    Derajat anonimitas (\textit{degree of anonymity}) yang menormalkan nilai entropi terhadap batas atas teoritis sehingga berada pada rentang $[0,1]$ &
    \pageref{first-degree-symbol} \\ \hline

    \textbf{$H(P)$} &
    Entropi Shannon dari distribusi probabilitas $P$ atas himpunan anonim yang digunakan sebagai ukuran derajat ketidakpastian &
    \pageref{first-entropy-symbol} \\ \hline
\end{xltabular}

} % end todo